\section{The graphical group and its embedded sheets}
\label{sec:geometry}

We now realize the sampled voltage cover inside a Cayley graph.  A private
completion first makes the cover connected.  We then impose its closed-walk
words as group relators and prove that the cover embeds in the resulting
Cayley graph and its Heisenberg deck group embeds in the group.  The
translated copies of the cover
are the sheets used in the rest of the proof.  Our second goal is to control
polygons through distinct sheets while keeping their contact vertices fixed;
this will give the intersection bounds needed in Section~\ref{sec:domains}.

The argument uses graphical small-cancellation methods developed by
Gromov, Ollivier, and Gruber
\cite[\S2.3]{Gro03}\cite[\S\S2--3]{Oll06}
\cite[Definition~1.3 and Lemma~4.1]{Gru15}.
In particular, the reductions below compare lifted paths modulo
label-preserving deck transformations and then apply a finite-diagram
curvature argument.  We give the required form with its varying scale,
bridges, and prescribed boundary contacts.  The graphs in this section may
be infinite; every path, diagram, and word used in a proof is finite.

\subsection{Conventions and the input from the preceding construction}

An edge of a labeled graph has a specified positive orientation and one
formal letter.  Traversing it backwards reads the inverse formal letter.
A labeling is \emph{folded} if, at each vertex, each signed formal letter is
read on at most one outgoing edge incidence.  In particular, a path and its
initial vertex are determined by its word whenever that path exists.  Word
products are read in traversal order.  Thus the formal right Cayley graph
of a group with alphabet $S$ has one positive edge
\[
 (g,s)\colon g\longrightarrow gs\qquad(g\in\Gamma, s\in S).
\]
These are distinct cells indexed by $(g,s)$, even if their unordered endpoints
coincide with those of another cell.  We never replace this graph by an
undirected simple graph.  Left translation preserves all its labeled,
oriented cells.  The corresponding right regular representation is
$R(a)\delta_g=\delta_{ga^{-1}}$; the unitary
$I\delta_g=\delta_{g^{-1}}$ satisfies $IR(a)I=\lambda(a)$.

Write $X_{\mathrm o}$ for the surviving old graph, with its old alphabet
$S_{\mathrm o}$ and chosen voltages $h_e\in\HH$ on its positive edges;
put $h_{\bar e}=h_e^{-1}$.  Here and throughout this Section
\[
 a=\frac b{\log_2 q}\le 10^{-6},\qquad
 L=1.003,\qquad m_{\min}\ge10^6.
\]
The use of $a$ in this Section is solely for the slope of the scale function.
At an old vertex of level $i$ put $m(v)=m_i$.
The inputs furnished by Propositions~\ref{prop:tower-labels},
\ref{prop:pruning}, and~\ref{prop:sampled-symbols}, in the form used here, are:
\begin{enumerate}
\item $X_{\mathrm o}$ is countable, folded, and locally finite, with maximum
 degree at most $K_{b,\epsilon}D$.  Its edges join the same or consecutive
 levels, so $|m(v)-m(w)|\le a$ on every old edge.
\item If an immersed old path starts at $v$ and its word has two different
 projected occurrences in $X_{\mathrm o}$, its length is less than $Lm(v)$.
 The assertion applies to either occurrence and to the reversed path.
\item In the old voltage graph, every simple cycle $C$ satisfies
 $|C|>6.03m(v)$ at every vertex $v$ of $C$.
\end{enumerate}
For the second assertion, the recovery length from the preceding construction
is $\lceil1.002m(v)\rceil$; the bound displayed here follows because
$\lceil1.002m(v)\rceil<1.003m(v)$ when $m(v)\ge10^6$.
An immersed path means one with no immediate edge reversal.  Foldedness
identifies this condition with free reduction of its word.
These three explicit assertions, rather than a small-cancellation theorem,
are the inputs to the proofs below.

\subsection{A private completion with a controlled scale}

The following construction uses group-valued voltage covers in the sense
of Gross and Gross--Tucker \cite{Gross74,GT77}; we specify the multiplication
convention to fix the later operator formulas.
The added paths connect the old components, and two added loops make the
cover connected.  Letters used on single added base edges will also identify the full
label-preserving automorphism group with the deck group.  Their lengths
are chosen to preserve the local girth and recovery bounds.

\begin{lemma}[Private completion]\label{lem:private-completion}
There is a connected, countable, folded labeled graph $X$ containing
$X_{\mathrm o}$ as its old subgraph, and an extension of the voltages, with
the following properties.
\begin{enumerate}
\item Each added edge has its own new positive letter, used on no other
 base edge.  Every new vertex has degree two and has only private incident
 edges.  The maximum degree is at most $K_{b,\epsilon}D+6$.
\item The function $m$ extends to the new vertices, remains at least
 $m_{\min}$, and satisfies $|m(v)-m(w)|\le a$ on every edge.
\item The voltage cover $\widetilde X$, with edges
 \begin{equation}\label{eq:voltage-convention}
 (v,h)\longrightarrow(w,hh_e)
 \end{equation}
 over $e\colon v\to w$, is connected.  Its deck action is the left action
 $f(v,h)=(v,fh)$ of $\HH$.
\item Every simple cycle $C$ of $\widetilde X$ satisfies
 $|C|>6.03m(v)$ at every vertex $v$ of $C$.
\item If two paths in $\widetilde X$ have the same word and their projected
 occurrences in $X$ differ, neither contains a private edge.  If they are
 immersed and start at $v,w$, their common length is less than both
 $Lm(v)$ and $Lm(w)$.
\end{enumerate}
The scale on the cover is the pullback of the scale on $X$.
\end{lemma}

\begin{proof}
Choose an old vertex in each component of $X_{\mathrm o}$ and enumerate the
components.  Join the chosen vertices in a chain by new subdivided paths;
if there are finitely many components use a finite chain.  The chain has
no new branching vertices.  For each joining path with endpoints $v,w$,
choose its integer length $n$ to satisfy
\begin{equation}\label{eq:private-length}
 n\ge\max\left\{20\max\{m(v),m(w)\},\ \frac{|m(v)-m(w)|}{a}\right\}.
\end{equation}
Interpolate $m$ linearly along the path.  All of its voltages may be the
identity.  At one old basepoint $v_*$ attach two further subdivided based
loops, each of length at least $20m(v_*)$, with new interior vertices.
Keep $m=m(v_*)$ on these loops.  Assign their voltages so that their ordered
products are respectively $x$ and $y$; for example use identity voltages
except on one positively traversed edge of each loop.  Give every added
positive edge a distinct new letter.  This preserves foldedness.  The chain
adds at most two incidences at any old vertex and the two loops add four
at $v_*$.  Their lengths are greater than two, so all new vertices have
degree two.  The interpolation proves the second assertion.

The completed base graph is connected.  A path from $v_*$ to a given base
vertex $v$ has some voltage $h_0$.  By first traversing based words in the
two added loops, one can obtain any prescribed voltage in $\HH$, since
$x,y$ generate $\HH$.  Following the chosen path thereafter reaches any
specified lift $(v,h)$.  Thus the cover in
\eqref{eq:voltage-convention} is connected.  Its displayed left action is
free on vertices and transitive on each fiber.

Consider a simple cycle that uses a private edge.  Since every private
interior vertex has degree two, the cycle traverses an entire lift of a
joining path or one of the two added subdivided loops.  Let $v$ be an old
endpoint of that lifted path and let $\ell$ be the length of the cycle.
Then $\ell\ge20m(v)$.  At any other vertex $w$ of the cycle the Lipschitz
bound gives $m(w)\le m(v)+a\ell$.  Consequently
\[
 6.03m(w)\le6.03(1/20+a)\ell
 \le0.30150603\ell<\ell.
\]
A cycle using only old edges satisfies the required inequality by the input
hypothesis.  This proves the fourth assertion, and in particular the cover
has neither loops nor two-edge circuits.

A private letter in a common word identifies its unique oriented base edge.
Foldedness then propagates equality of the two projected paths forwards
and backwards from that edge, without any assumption that the paths are
immersed.  This proves the first statement of the last assertion.
When different projected occurrences are immersed they therefore consist
of old edges, and the old recovery bound applies.  Reversing the paths
gives the bounds at their other endpoints as well.
\end{proof}

For future use, every nonempty cyclically immersed closed walk $W$ based
at $v$, including a nonsimple one, has
\begin{equation}\label{eq:closed-walk-lower}
 |W|>\frac{6.03}{1+6.03a}\,m(v).
\end{equation}
Indeed it contains a simple cycle based at some vertex $w$ reached within
$|W|$ steps.  Thus
$|W|>6.03m(w)\ge6.03(m(v)-a|W|)$.
Also, if consecutive immersed comparison paths, each bounded by $Lm$ at
its own initial point, number at most $j$, their total length is at most
\begin{equation}\label{eq:arc-geometric-sum}
 L\sum_{t=0}^{j-1}(1+La)^t\,m(v).
\end{equation}
This follows by induction from the Lipschitz bound.
In particular,
\begin{align}
 L\sum_{t=0}^2(1+La)^t&<3.009004<3.014,
 \label{eq:three-arcs}\\
 L\sum_{t=0}^5(1+La)^t&<6.018016
 <6.029963<\frac{6.03}{1+6.03a}.
 \label{eq:six-arcs}
\end{align}

\subsection{Finite word diagrams and their perimeter budget}

Put $S=S_{\mathrm o}\sqcup S_{\mathrm p}$, where $S_{\mathrm p}$ is the
countable private alphabet.  Write $\operatorname{lab}(C)$ for the word
read along an oriented path $C$, and define
\begin{equation}\label{eq:graphical-presentation}
 \Gamma=\left\langle S\ \middle|\
 \operatorname{lab}(C)=1\text{ for every finite closed walk }C\text{ in }\widetilde X
 \right\rangle.
\end{equation}
Each occurrence of a relator in a diagram will retain an actual closed
walk of $\widetilde X$ reading it.  Two occurrences of the same free word
need not be the same projected walk.

The relators make the group element represented by a path from a fixed
source vertex depend only on its endpoint.  To obtain an embedding, we
must prove the converse: a source path with trivial label has equal
endpoints.  We will use diagrams
to prove this closure property.  Keeping the source walks lets us distinguish
two genuinely different projected occurrences from two lifts of the same
base path.

We use finite plane multigraph diagrams.  Such a diagram is a finite
connected graph embedded in the sphere, with one complementary face
specified as the exterior.  Each edge is oriented and labeled by a formal
letter.  Every other face has a specified source closed walk in
$\widetilde X$ reading its oriented facial contour.  A facial contour is a
walk and can repeat vertices and bridges.  Perimeters count edge
incidences, including both traversals of a bridge.  The \emph{filling cost}
is the sum of the perimeters of the bounded faces.  A one-vertex graph
with no edges is allowed.

\begin{lemma}[Elementary word-diagram construction]\label{lem:word-diagram}
If a finite word $w$ is trivial in \eqref{eq:graphical-presentation}, there
is a finite plane multigraph diagram with exterior contour reading the
literal word $w$.  More precisely, if in the free group on $S$
\[
 w=\prod_{j=1}^n u_jr_ju_j^{-1},
\]
where each $r_j$ is read on a chosen closed source walk, a diagram may be
chosen with filling cost at most $\sum_j|r_j|$.
Conversely, the exterior word of any such diagram is a product of
conjugates of its bounded-face words in the free group, using each
bounded facial word once with the compatible orientation.  Its labeled
one-skeleton therefore has a well-defined vertex map to the formal Cayley
graph after prescribing the image of one vertex.
\end{lemma}

\begin{proof}
Only the finitely many letters in the displayed expression matter.  Let
$B$ be their finite bouquet of oriented circles.  The elementary
free-cancellation rule constructs a null homotopy in $B$ for a freely
trivial word: delete adjacent inverse traversals one pair at a time.
Take an oriented sphere with one disk removed for each of the words $w$
and $r_1,\ldots,r_n$, reading $w$ with the opposite boundary orientation.
Choose a cutting tree joining these boundary circles, and prescribe the
words $u_j$ on its cuts.  Cutting along the tree gives a disk whose
boundary reads the freely trivial word supplied by the displayed
identity.  The cancellation homotopy extends its boundary map to $B$.
Regluing the cuts gives a map from the punctured sphere to $B$ with the
specified word maps on all boundary circles.

Choose one interior point of each circle of $B$.  A piecewise linear
approximation relative to the boundary can be made transverse to these
points as follows: triangulate the source sufficiently finely so each
triangle maps into an interval of a circle or into a small tree around
the bouquet vertex; choose the interior points outside the images of
source vertices.  Inside each triangle the inverse image of a selected
point is a line segment or is empty.  These segments fit to disjoint arcs
and circles.  Their boundary endpoints are precisely the occurrences of
the selected letters on the word boundaries, and each arc joins opposite
signed occurrences of the same letter.  Closed inverse-image circles may
be ignored.  Thicken each remaining arc to a narrow disjoint band, cap
the word boundaries by disks, and contract each cap to a vertex and each
band to an edge.  This is an embedded ribbon multigraph in the sphere;
its cyclic order at a word vertex is exactly the cyclic order of that
word's letters.

If $w$ is nonempty, retain the component containing its word vertex and
take the planar dual, specifying that vertex's dual face as exterior.
Every bounded dual face corresponds to one of the retained $r_j$ word
vertices, with its original source walk, so the perimeter budget follows.
This argument permits loops, multiple edges, and nonsimple dual face
contours.  If $w$ is empty, the one-vertex empty diagram suffices.  A
freely trivial but nonempty $w$ is covered by the same construction, even
when there are no relator word vertices in its component.

For the converse, take a thin regular neighborhood of the embedded
graph.  Its inner boundary circles correspond to the bounded facial
contours and its outer boundary to the exterior contour.  Cut between
these boundary components along a tree of arcs.  The cut surface is a
disk, and the neighborhood retracts to the labeled graph.  Reading the
disk boundary in $B$ gives the exterior word as a product of conjugates
of the bounded face words, each occurring once with its compatible
orientation.  The same conclusion holds for the contour of each simple
cycle of the graph, by taking the portion on its bounded side.
Consequently every closed walk in the graph is trivial in $\Gamma$:
split a walk at repeated vertices until it consists of simple cycles
and immediate edge reversals.  Labels of paths from a chosen vertex to
any other vertex are therefore independent of the chosen path in
$\Gamma$.  This is the asserted vertex map; each edge maps to its formal
Cayley cell.
\end{proof}

Two elementary facts about plane multigraphs will be used repeatedly.
An edge has the same complementary face on both sides if and only if it
is a bridge.  Removing such an edge separates the graph into two
components.  Between its two traversals, the incident face follows the
exterior contour of the appropriate component.  These assertions follow
by the simple-cycle criterion for a bridge and the Jordan separation
property of a simple cycle; they remain valid when a contour repeats
other bridges.

\subsection{Null words close in the cover}

\begin{lemma}[Closure of a null-labeled source path]\label{lem:null-path-closes}
A path in $\widetilde X$ whose label is trivial in $\Gamma$ has equal
initial and terminal vertices.
\end{lemma}

\begin{proof}
We argue by a minimal counterexample.  Coincident projected edge occurrences
will permit reductions that lower the total face perimeter.  The remaining
comparisons are short by recovery, so each bounded face needs at least seven
arcs; the final Euler count will rule out such a diagram.

Suppose a counterexample exists.  Among all such source paths and all
finite diagrams for their words, minimize first the filling cost $c$,
then the exterior path length, and finally the number of diagram edges.
All three quantities are nonnegative integers.  The exterior path has
one marked end-to-start join.  Its source endpoints are required to be
distinct, but may change when considering a competing counterexample.
This global choice is needed for the bridge argument below.

Each bounded source face may be assumed nonempty and cyclically
immersed.  To justify this without assuming anything about diagram
regularity, use the free-group face expression from
Lemma~\ref{lem:word-diagram}.  Foldedness turns every free cancellation in
a source walk into an actual edge reversal.  A freely trivial closed
face word can be omitted from that expression.  Otherwise delete its
backtracks and its cyclic conjugating tail, absorbing the tail into its
conjugator in the expression.  The remaining word is read on a closed
source walk with smaller perimeter.  The construction part of
Lemma~\ref{lem:word-diagram} gives a diagram for the unchanged exterior
word with strictly smaller filling cost.  Minimality excludes these
possibilities.  In the same way the exterior source path is linearly
immersed: deleting any backtrack preserves its two source endpoints,
does not increase the least filling cost, and decreases its length.

We next compare the two source occurrences of each diagram edge incident
to a bounded face, orienting them to read the same signed letter.  There
are three cases, according to the two face incidences.  In each case,
coincident projections can be eliminated; a bridge on one bounded face
can be eliminated without even requiring coincident projections.

First, if two different bounded faces project their common edge to the
same base edge of $X$, a unique deck transformation aligns the two
lifted edges.  Delete this diagram edge.  It is not a bridge, so the
remaining graph is connected and the faces merge.  Splice their aligned
closed source walks, omitting the two opposite edge traversals.  The
new source walk is closed, and its perimeter is the sum of the old two
perimeters minus two.  Other shared edges are retained and may occur
twice on the new contour.  The new diagram is permitted and has smaller
filling cost, a contradiction.

Second, suppose an edge has the same bounded face $F$ on both sides.
It is a bridge.  The exterior face of the original diagram, including
its marked join, lies entirely on one side of that bridge.  Let $D_1$
be the other component after deleting the bridge.  In the source walk
assigned to $F$, the excursion into $D_1$ is a subpath $\chi$ reading the
bridge, the exterior contour of $D_1$, and the reverse bridge.  All
bounded faces of $D_1$ are former bounded faces other than $F$.
Lemma~\ref{lem:word-diagram} shows that $\operatorname{lab}(\chi)=1$ in $\Gamma$ and
gives a filling of cost strictly less than $c$, because the positive
perimeter of $F$ is absent.  If the source endpoints of $\chi$ differed,
it would be a smaller counterexample to the globally chosen one.
They therefore agree.  Delete this source excursion, the bridge, and
$D_1$.  The remaining walk assigned to $F$ still closes, the original
exterior source path is unchanged, and the filling cost has decreased
by at least the two bridge incidences.  This is again impossible.
Notice that this argument did not assume embeddedness and did not need
the two projected bridge occurrences to coincide.

Third, if a bounded face and the exterior path project their common edge
to the same edge of $X$, align their lifts by a deck transformation.
Replace this step of the exterior source path by the complementary path
of the bounded face, oriented to have the same lifted endpoints.  Delete
the shared diagram edge, thereby absorbing that bounded face into the
exterior.  The source endpoints of the exterior path remain distinct,
and its filling cost drops.  Any subsequent source backtrack deletions
preserve that failure.  This also contradicts minimality.

These reductions leave only different projected occurrences along edges
incident to bounded faces.  We now collect them into arcs.  Retain the
marked join and suppress all other vertices of degree two.
There is no unmarked vertex of degree one: it would force an immediate
reversal in a cyclically immersed bounded face or within the linearly
immersed exterior path.  The marked vertex has degree at least one if
any edge is present.  Each resulting edge, called an \emph{arc}, has
constant face incidences and a constant exterior-path incidence along
its interior.  An arc incident to a bounded face reads an immersed word
on both its occurrences.  By the preceding reductions those projected
occurrences differ.  Lemma~\ref{lem:private-completion} therefore bounds
its length by $Lm$ at the initial point of either occurrence.

The length bounds force every bounded face to have at least seven arc
incidences.  Otherwise apply
\eqref{eq:arc-geometric-sum} to its at most six arcs.  Its nonempty
cyclically immersed closed source walk would violate
\eqref{eq:closed-walk-lower} and~\eqref{eq:six-arcs}.

If there are no bounded faces, the diagram is a tree.  A nonempty finite
tree has at least two leaves, whereas only the marked vertex could be
a leaf.  The empty diagram has empty word and cannot represent a
counterexample.  Thus at least one bounded face is present.  Suppression
is also defined when the graph is a lone cycle, because its marked join
is retained.  Write $V,E,F$ for the compressed vertex, edge, and total
face counts.  Euler's formula and the vertex degrees give
\begin{equation}\label{eq:exceptional-curvature}
 6F-2E
 =12+2\sum_v(\deg(v)-3)\ge8.
\end{equation}
On the other hand this expression equals
$\sum_f(6-\deg(f))$.  Each bounded-face term is at most $-1$, and the
one exterior term is at most $6$.  Their sum is less than $8$, a
contradiction.  The retained exceptional join is the reason for using
$8$, rather than assuming the stronger bound $12$.
\end{proof}

The closure lemma supplies the injectivity needed to pass from the cover
to its Cayley image.  The private letters then determine which translations
preserve that image.

\begin{proposition}[Graphical realization]\label{prop:graphical-realization}
There is an injective map of vertices and formal oriented edges
\[
 j\colon\widetilde X\longrightarrow\operatorname{Cay}(\Gamma,S)
\]
preserving labels and sending $(v_*,1)$ to the identity.  The subgroup
of vertices over $v_*$ defines an injective homomorphism
$\phi\colon\HH\hookrightarrow\Gamma$.  If
$t_v=j(v,1)$, then
\begin{equation}\label{eq:deck-transport}
 j(v,h)=\phi(h)t_v.
\end{equation}
The full group of label-preserving automorphisms of $\widetilde X$ is its
deck group.  Consequently the distinct translated images, called
\emph{sheets}, are indexed by $\Gamma/\phi(\HH)$, and the stabilizer of
$g j(\widetilde X)$ under left translation is
$g\phi(\HH)g^{-1}$.
\end{proposition}

\begin{proof}
Assign to a source vertex the element represented by any path from
$(v_*,1)$ to that vertex.  Two choices differ by a closed source walk,
whose word is a relator, so the assignment is well defined.  If two
vertices have the same image, a source path between them has null label;
Lemma~\ref{lem:null-path-closes} identifies its endpoints.  This proves
vertex injectivity.  Edges with the same formal Cayley image have the
same initial image and the same signed label.  Their initial vertices
agree by injectivity and their source edges agree by foldedness.  Hence
formal edges also embed.

For $f\in\HH$, let $\phi(f)=j(v_*,f)$.  A source path from $(v_*,1)$
to $(v_*,f_1)$, followed by the left deck translate by $f_1$ of a path
to $(v_*,f_2)$, ends at $(v_*,f_1f_2)$ and reads the product of their
words.  Thus $\phi(f_1f_2)=\phi(f_1)\phi(f_2)$.  It is injective by
vertex injectivity.  Following a path to $(v_*,h)$ and then the left deck
translate of a path to $(v,1)$ proves \eqref{eq:deck-transport}.
More generally, connectedness and label preservation show
\[
 j(fu)=\phi(f)j(u)\qquad(u\in\widetilde X).
\]

Fix one of the private positive base edges.  Its letter appears on no
other base edge.  A label-preserving source automorphism must send a
chosen lift of this edge to a lift of the same base edge.  Compose the
automorphism with a deck transformation to fix that edge's initial
vertex.  An automorphism of a connected folded labeled graph fixing a
vertex fixes every path from that vertex, hence every vertex and edge.
It is the identity.  Thus every source automorphism is a deck
transformation.

If a left translation by $g$ stabilizes $j(\widetilde X)$, its
restriction, conjugated by $j$, is a label-preserving source
automorphism.  It is deck multiplication by some $f$.  Evaluating at
$(v_*,1)$ gives $g=\phi(f)$.  Conversely \eqref{eq:deck-transport} shows
that every $\phi(f)$ stabilizes the image.  The remaining stabilizer and
coset assertions follow by conjugation and translation.
\end{proof}

We identify $\HH$ with its image from now on.  Each sheet is equipped
with its source vertices, edges, and scale $m_Q$.  These do not depend on
the representative $g$ of its coset: changing $g$ by an element of $\HH$
changes the source parametrization by a deck transformation, which
preserves the base vertex and scale.  There is a useful rigidity
observation.  If two source-parametrized sheet maps agree at one source
vertex after deck alignment, they agree everywhere by connectedness and
label propagation.  Thus two different sheets meeting at a physical
vertex cannot have the same projected base vertex there.  In particular,
a common path in two different sheets always has different projected
occurrences.  This observation concerns nonconstant paths; it does not
exclude isolated contacts at private vertices.

\subsection{Polygons with fixed sheet contacts}

We next control polygons whose vertices are prescribed intersections of
sheets.  The sides may change during diagram reduction, but their contact
vertices must stay fixed.  This is the form needed to compare the intrinsic
distances between contacts in Section~\ref{sec:domains}.

\begin{lemma}[The exterior-arc bound]\label{lem:sheet-polygon}
Let $Q_1,\ldots,Q_k$ be distinct sheets, $k\ge2$, and let
$g_i\in Q_i\cap Q_{i+1}$, with indices modulo $k$.  Contacts may
coincide.  Choose a source path in $Q_i$ from $g_{i-1}$ to $g_i$ for
each $i$.  There is a finite diagram and a choice of these paths,
\emph{with every prescribed endpoint $g_i$ fixed}, such that:
\begin{enumerate}
\item Each side is linearly immersed.  After retaining all junction
 vertices and suppressing other degree-two vertices, every arc has two
 different projected occurrences in $X$.  Each occurrence is immersed,
 contains no private edge, and has length less than $Lm_Q(v)$ at its
 initial point.  Each entire nonconstant side is the concatenation of
 its exterior arc occurrences.
\item Every bounded face has at least seven arc incidences.
\item If $B$ counts exterior arc incidences, counting an exterior bridge
 twice, then
 \begin{equation}\label{eq:polygon-exterior}
 B\le4k-6.
 \end{equation}
\end{enumerate}
A diagram with no edges has $B=0$ and all its sides constant.
In particular, a nonconstant side cannot have a private vertex as an
endpoint.  A side whose intrinsic endpoint distance is greater than
$3.014m_Q(v)$ at its initial endpoint has at least four arc incidences.
\end{lemma}

\begin{proof}
The concatenation of any chosen sheet paths is a closed walk in the
Cayley graph and has null word.  Lemma~\ref{lem:word-diagram} supplies a
finite diagram.  Over all choices with the prescribed contacts fixed,
minimize lexicographically the filling cost, the total side length, and
the number of edges.  Retain every boundary vertex carrying a junction
mark; there are at most $k$ such vertices, even if several junctions
coincide.

The source-path and bounded-face reductions at the beginning of
Lemma~\ref{lem:null-path-closes} apply unchanged: bounded source faces
are nonempty and cyclically immersed, while each side is linearly
immersed.  Coincident projected edges between two bounded faces can be
merged after deck alignment, decreasing filling cost.  A bridge with
both incidences on one bounded face can be removed as in that proof.
Here there is no preliminary injectivity issue: its source excursion
has null label by the cut-off component, and
Lemma~\ref{lem:null-path-closes} makes it close.  The component containing
the exterior retains all marked contacts.

If a bounded face and a boundary side have coincident projected
occurrences along an edge, align the closed face lift to that side and
replace the one side step by the complementary face path.  The two
source endpoints of the replaced step agree with the old ones, and
therefore both prescribed endpoints of the full side remain fixed.
Deleting the separating edge absorbs the face and decreases filling
cost.  This is impossible by minimality.

It remains to analyze edges whose two sides are exterior.  Such an edge
is a bridge.  If its two traversals occur in the same side, the interval
between them within that side is an excursion with no junction in its
interior.  Its endpoints have the same Cayley vertex, and hence the
same source vertex by Proposition~\ref{prop:graphical-realization}.
Delete the excursion from that side and remove the bridge and its
cut-off component.  No prescribed junction is lost.  Filling cost does
not increase and total side length decreases, again contradicting
minimality.  If instead the two traversals belong to different sides,
their sheets are distinct.  Coincident projected occurrences would be
deck-aligned and agree at a physical vertex; the rigidity observation
after Proposition~\ref{prop:graphical-realization} would identify the
two sheets.  Thus these projected occurrences differ too.

Suppress every unmarked degree-two vertex.  An unmarked leaf would
force a backtrack in one side or one bounded face, so all unmarked
vertices have degree at least three.  Marked vertices have degree at
least one when edges occur.  The face and side incidences of each
resulting arc are constant along its interior: a side change occurs at
a retained mark, and a face change requires branching.  Both of its
occurrences therefore lie within immersed source paths.  The preceding
reductions show they have different projections.  Apply
Lemma~\ref{lem:private-completion} to obtain all the assertions about
arc lengths and private edges.  The proof using
\eqref{eq:closed-walk-lower} and~\eqref{eq:six-arcs} again gives at least
seven arcs on every bounded face.

For completeness the Euler count uses the exterior degree itself,
rather than the coarser estimate by all edges.  Let $f$ be the number
of bounded faces and let $V,E$ count the compressed vertices and edges.
When edges occur,
\[
 E=V+f-1,\qquad 2E\ge3V-2k,\qquad E\le3f+2k-3.
\]
Every edge has two face incidences, so
\[
 2E-B=\sum_{\text{bounded }F}\deg(F)\ge6f.
\]
Consequently $B\le2E-6f\le4k-6$.  If the diagram has no edges, all
of its marked vertices coincide and every side is constant, giving
the stated separate case.

Every edge incident to a private source vertex is private.  If a side
with such an endpoint were nonconstant, its initial or terminal arc
would contain a private edge, contradicting the arc conclusion already
proved.  Finally, \eqref{eq:three-arcs} bounds a side of at most three
arcs by $3.014m_Q(v)$ at its initial endpoint.  The contraposition proves
the last assertion.
\end{proof}
